\documentclass[../main.tex]{subfiles}
\begin{document}

\chapter{Mathematical modelling of atmospheric flows}
\label{chap:MODELLING}

\begin{compactenum}
	\item section \Cref{sec:IBCONDITIONS} discusses suitable boundary conditions the system needs to be accompanied by;
	\item section \Cref{sec:WELL_BALANCING} is devoted to the \emph{well-balanced} property of certain discretisation schemes.
\end{compactenum}

\section{Initial and boundary conditions}
\label{sec:IBCONDITIONS}

\subsection{Radiation condition}
\label{sec:RADIATION}

\subsection{Terrain-following coordinates}
\label{sec:SIGMA}

\subsection{Inflow and outflow boundaries}
\label{sec:INFLOW-OUTFLOW}

\section{Well-balanced methods}
\label{sec:WELL-BALANCED}

This section provides an introduction into specialized numerical methods for the compressible Euler equations. The motivation for the deployment of such methods is presented in the first subsection, whereas our implementation is described in \Cref{sec:WELL-BALANCED-SPH}. Results obtained with the use of these methods are then presented in \Cref{chap:EXPERIMENTS}.

\subsection{Preservation of the hydrostatic balance}
\label{sec:INEXACT}

Internal gravity waves can be considered perturbations of the background atmospheric state. It is standard to take the background atmospheric state as \emph{hydrostatic}\sidenote{As we did in the section \Cref{sec:MTN-WVS}.}: $\vb{v} = 0$ and the pressure $p$ and density $\rho$ solve the equation \Cref{eq:HYDROSTATIC-P}

\begin{equations}[single,*]
\grad p = \rho \vb{g}.
\end{equations}

In \Cref{sec:HYDROSTATICS} we have shown that for an \emph{isothermal} atmosphere, the resulting pressure and density read \Cref{eq:HYDROSTATIC}

\begin{equations}[columns,*]
	\rho(z) &= \rho_0\exp\left(- \frac{z}{H_s}\right), \\
	p\left(z\right) &= p_0\exp\left(- \frac{z}{H_s}\right). 
\end{equations}

Suppose now we wish to \emph{numerically} solve the \emph{non-stationary} isentropic compressible Euler system with the initial condition given as the hydrostatic balance, \ie solve\sidenote{We are using a formulation of the Euler equations that is more common in numerical mathematics.}

\begin{equations}[columns,!,numberline=all]
\label{eq:NONSTAT-EULER}
\partial_t \rho + \div\left(\rho \vb{v}\right) &= 0, \\
\partial_t\left(\rho \vb{v}\right) + \div\left(\rho \vb{v}\right) &= - \grad p + \rho \vb{g}, \\
\partial_t \theta + \div\left(\theta \vb{v}\right) &= 0,
\vb{v}\left(0, \vb{x}\right) = \vb{0}, \theta\left(0, \vb{v}\right) &= \theta_0\left(\vb{x}\right), \rho\left(0, \vb{x}\right) = \rho_0\exp\left(- \frac{z}{H_s}\right),
\end{equations}

together with $p = p\left(\theta, \rho\right)$ given and some boundary conditions.\sidenote{A rigorous numerical treatment of the boundary conditions is complicated, see \eg \autocite[p.~59]{levequeFiniteVolumeMethods2002}.} \emph{Mathematically} speaking, since at time $t=0$ the system is in hydrostatic equilibrium, the right-hand side of the momentum equation as well as the convective term vanish and the system \emph{remains} in the hydrostatic equilibrium. Yet, in a numerical simulation, generally the hydrostatic equilibrium is \emph{not preserved},

\begin{equations}[single,*]
\grad p_{i} - \rho_i \vb{g} \neq 0,
\end{equations}

where $p_i, \rho_i$ are some approximations of the solution \enquote{at time $i$} obtained by the numerical method. Since the right-hand side of \Cref{eq:NONSTAT-EULER} does not vanish on the \emph{discrete level}, a spurious non-physical force is introduced and the gas moves with \emph{nonzero velocity.}

\subsection{Residual forces \& Waves}
\label{sec:WB-WAVES}

The failure to preserve hydrostatic balance is particularly damaging for simulations of internal gravity waves. As mentioned throughout the text, internal gravity waves are small perturbations of the \emph{background state.} When the hydrostatic balance is inexact, the spurious forces influence the simulation and affect the flow field. 

In \autocite[p.~6]{kappeliWellbalancedFiniteVolume2016a} it is shown that, with the used numerical method, for small perturbations of the velocity field\sidenote{The authors used a sine perturbation of the form \begin{equations}[single,*]
\delta v = \sin\left(\omega t\right),
\end{equations}
with A = \num{10e-8}.}
the effect of the spurious forces is considerable. For large-amplitude flows, in \autocite[p.~6]{kappeliWellbalancedFiniteVolume2016a} it is shown that the effect is not critical. At least in the case of small-amplitude wave dynamics, the preservation of hydrostatic balance remains an important property. 

\subsection{Well-balanced schemes}
\label{sec:WB-SCHEMES}

Numerical methods that preserve the \emph{analytic} steady solutions on the \emph{discrete} level are called \emph{well-balanced}\autocite{cargoSchemaEquilibreAdapte1994,greenbergWellBalancedSchemeNumerical1996}. The methods had been first used for shallow water equations \autocite{greenbergWellBalancedSchemeNumerical1996} and is nowadays standard in many areas where numerical solution of \emph{hyperbolic} PDEs is required, \eg atmospheric flows\autocite{kappeliWellbalancedFiniteVolume2016a,chertockWellbalancedNumericalMethod2023,berthonFullyWellbalancedEntropystable2025} and astrophysics \autocite{kappeliWellbalancedMethodsComputational2022}. 
In those applications, \emph{grid methods} like FVM or DG are used. In \Cref{sec:WELL-BALANCED-SPH} we derive \emph{well-balanced} formulations of SPH suitable for an isentropic atmospheric flow.




\end{document}


